Why is it hard to guarantee data durability in a decentralized storage system?
In this section, we analyze key challenges faced by existing systems to prevent data losses in an open, permissionless environment.
We then present our approach of using rateless fountain code and verifiable randomness to tackle the durability challenge.

\subsection{The Durability Challenge in Decentralized Storage}
\label{sec:approach:challenge}

How to maintain data durability in the presence of server and device failures?
The de facto standard approach is data replication and erasure code.
It is well-studied how \textit{redundancy} can mask failures -- including Byzantine ones -- of individual components and ensure data availability.
The approach is particularly effective in centralized settings, in which system participants and failure distributions are managed by a single entity.
For instance, the central administrator can control node/device selection such that out of the $n$ nodes storing a data object, at most $n-r$, where $r$ is the level of redundancy, of them can fail.
\lijl{Can add placement group discussion here.}

The formula changes completely in an open, decentralized system.
Nodes are free to join and leave the system at any time.
There is no centralized control over who are the participants and their trustworthiness.
In these systems, adversarial behaviors are the \textit{common case}, not the exception.
Adversaries have strong attack power.
They can partition regions of the network using BGP poisoning or distributed denial-of-service (DDoS);
they can collude with each other and compromise non-faulty nodes;
they can create and change identities in a highly dynamic fashion.

A direct implication of a decentralized environment is that failure assumptions made by prior systems no longer hold.
Adversaries can launch Sybil attacks such that they control more than $n-r$ copies or erasure code symbols of an object;
once they locate more than $n-r$ nodes responsible for an object, they can also launch targeted attacks to partition their network or compromise their systems.
Data durability is immediately violated under these attack scenarios.

Existing blockchain systems face similar attack vectors.
However, decentralized storage systems are particularly vulnerable, given their scale and redundancy constraints.
In large scale deployments, the size of the system, which can reach beyond 100K nodes~\cite{ipfs}, easily overwhelms the replication factor of each storage object, which is typically below 20 \lijl{citation?}.
In fact, reducing redundancy is a key metric to improve the overall storage efficiency.
This is in contrast to a blockchain, in which chain data is replicated on every participant.
With such a small redundancy, effectiveness of adversarial attacks grows exponentially, even when the adversaries account for a small percentage of the entire system.

Another challenge faced by open, decentralized systems is that consensus and coordination are \textit{costly}.
To safely confirm a transaction, the Bitcoin network requires a confirmation latency of around one hour;
even the faster Ethereum network incurs a finality latency of over 12 seconds.
Besides latency, both networks also suffers low consensus throughput -- 7 transaction per second (TPS) for Bitcoin and 20 TPS for Ethereum -- as well as high transaction fees.
Storage mechanisms that rely on frequent consensus are impractical in such settings.

\subsection{The Case for Using Fountain Code}

Fountain codes are a family of erasure codes that have two special properties.
Firstly, fountain codes are rateless: They generate a potentially \textit{infinite sequence} of encoding symbols from $k$ source blocks.
Any $k+\epsilon$ encoding symbols can be used to reconstruct the original source data.
This is different from traditional maximum distance separable (MDS) codes, such as Reed-Solomon codes, which produces $n$ fixed encoding symbols from the $k$ source blocks.
Second, fountain codes work well with large values of $n$, $k$, and $r$ --- real deployment of $(1200, 1000, 200)$ fountain code has been shown to be efficient~\cite{liquid}.
Such numbers are two orders of magnitude larger than practical MDS codes used in storage systems.

Why do these properties matter in a decentralized storage system?
Using traditional erasure codes with a small value of $r$, the storage system needs to quickly regenerate encoding symbols when just a few stored symbols are lost due to failures.
If the repair process is not done in a timely manner, the system risks losing the source data permanently.
However, decentralized systems, unlike centrally managed clusters, exhibit high degree of asynchrony: Network conditions vary both spatially and temporally; there exists large deviation in processing speed across nodes; transient unreachabilities are prevalent.
Strict timing requirement for detection and recovery of encoding symbols is therefore impractical in a decentralized environment.
On the other hand, large values of $r$ in fountain codes can tolerate bursty encoding failures without losing the data, and permits the system to recover at a steady average rate.
Such property is more favorable in a highly aynchronous environment.
Note that the increase in $r$ in fountain codes does not lead to higher redundancy: a \XXX and a \XXX code have the exact same redundancy, \XXX.

The rateless property of fountain code is also beneficial.
As discussed in \autoref{sec:approach:challenge}, coordination and consensus are costly in decentralized systems.
MDS codes with small values of $(n, k, r)$ necessitate participants to coordinate to agree on the assignment of symbols.
Uncoordinated symbol generation may lead to insufficient \textit{unique} symbols to recover the original data, an issue analogous to the coupon collector's problem \lijl{citation}.
Rateless codes eschew this issue completely.
The availability of an infinite stream of encoding symbols enables participants to \textit{independently} generate symbols with a negligible probability of collision.
\lijl{May need to rephrase this sentence.}


%Infinite, ordered fragments. Given order $i$, the $i$-th fragment can be generated in time linear to data object size.
%Similar arguments have been made in the liquid system~\cite{liquid} and OceanStore~\cite{oceanstore}

\subsection{Verifiable Randomness for Adversary Tolerance}

\textbf{Problem with replication/small code.} Fixed set of peers to be responsible, vulnerable to targeted attack.
Replication: high redundancy rate. Short repair interval, have to do eager repair.

\textbf{Problem with static large code.} Can only tolerance limited number of extra faulty nodes.
It's not practical to tolerance targeted attack only with erasure code (i.e. let $r > B$): either $k$ too large or too much redundancy.

\textbf{Problem with replication with rotation.} Predictable rotation is still vulnerable. Periodical downtime.

\textbf{Insight of randomized rotation.} Only encounter a small subset of Byzantine node for all the time.

\textbf{Overhead of fountain code.} Encode/decode performance is good enough. $\epsilon$ exist but small enough (on average much less than 1).
